\subsection{Characterizations of absolutely regular algebras}

PROPOSITION 8. Let \(k\) be a field and let A be an essentially of finite type \(k\)-algebra. Let I denote the kernel of the homomorphism of \(k\)-algebras \(\mu: A \otimes_k A \to A\) such that \(\mu(x \otimes y) = xy\) for \(x\) and \(y\) in A. The following conditions are equivalent:

\begin{enumerate}
\def\labelenumi{(\roman{enumi})}
\item
  the k-algebra A is absolutely regular;
\item
  for every regular k-algebra C, the ring A ⊗k C is regular;
\item
  the ring A ⊗k A is regular;
\item
  the ideal I of A ⊗k A is completely secant (§ 5, no. 1, def. 1).
\end{enumerate}

Let B denote the ring A \(\otimes_{k}\) A and endow it with the structure of an A-algebra deduced from the homomorphism \(\rho: \mathrm{A} \to \mathrm{A} \otimes_{k} \mathrm{A}\) such that \(\rho(x) = x \otimes 1\); then \(\mu\) is a homomorphism of A-algebras and induces, by passage to the quotient, an isomorphism from B/I onto A.

(i)⇒(ii): this follows from Proposition 7.

(ii)⇒(iii): it suffices to apply (ii) with C=k, then with C=A.

(iii)⇒(iv): the A-module B is free, hence faithfully flat. If the ring B is regular, then A is regular (§ 4, no. 5, prop. 8); the ideal I is then completely secant (§ 5, no. 3, prop. 2).

(iv)⇒(i): suppose the ideal I is completely secant and first prove that A is regular. Let m be a maximal ideal of A and let \(\nu:(\mathrm{A}/\mathfrak{m})\otimes_{k}\mathrm{A}\to\mathrm{A}/\mathfrak{m}\) be the homomorphism deduced from \(\mu\). The maximal ideal n=Ker \(\nu\) is equal to \(I((\mathrm{A}/\mathfrak{m})\otimes_{k}A)\); applying Proposition 6 of § 5, No. 6, to the A-algebra \(A'=A/\mathfrak m\), we see that the ideal n is completely secant in \((A/\mathfrak{m})\otimes_{k}A\). Consequently (§ 5, No. 3, Proposition 3), the local ring \(((\mathrm{A}/\mathfrak{m})\otimes_{k}A)_{n}\) is regular. Denote \(j:\mathrm{A}\to(\mathrm{A}/\mathfrak{m})\otimes_{k}\mathrm{A}\) the homomorphism \(x\mapsto1\otimes x\); since \(\nu\circ j\) is the canonical homomorphism from A to A/m, we have \(j^{-1}(\mathfrak{n})=\mathfrak{m}\). Thus j extends to a local homomorphism of local rings from \(A_{m}\) into \(((\mathrm{A}/\mathfrak{m})\otimes_{k}\mathrm{A})_{n}\), which makes the latter an \(A_{m}\)-module that is faithfully flat. By Proposition 8 of § 4, No. 5, the ring \(A_{m}\) is therefore regular. We have thus proved that A is regular.

Now let \(k'\) be an extension of \(k\). The kernel of the map \(\mu': A_{(k')} \otimes_{k'} A_{(k')} \to A_{(k')}\) deduced from the multiplication of \(A_{(k')}\) is none other than \(IA_{(k')}\); it is therefore completely secant in \(A_{(k')}\) (§ 5, No. 6, Proposition 6). The \(k'\)-algebra \(A_{(k')}\) therefore satisfies condition (iv); by what has just been seen, it is regular, and this proves (i).

Recall (A, III, pp. 133–134) that the quotient I/I², endowed with the A-module structure induced by ρ, is denoted Ωₖ(A) and called the module of k-differentials of A. When the k-algebra A is essentially of finite type, the ring A ⊗ₖ A is Noetherian, so that the A-module Ωₖ(A) is finitely generated.

One denotes by \(d_{A/k}\) or simply \(d\) the \(k\)-linear map from A into \(\Omega_k(A)\) which associates to an element \(x\) of A the class of \(x \otimes 1 - 1 \otimes x\) in \(\Omega_k(A)\). The map \(d\) is a \(k\)-derivation; for every A-module M and every \(k\)-derivation D: A → M, there exists a unique A-linear map \(g: \Omega_k(A) \to M\) such that D = g ∘ d (loc. cit., Proposition 18).

If S is a multiplicative subset of A, the canonical \(S^{-1}A\)-linear map (loc. cit., p. 136)

\[
\mathrm{S} ^ {- 1} \Omega_ {k} (\mathrm{A}) \rightarrow \Omega_ {k} (\mathrm{S} ^ {- 1} \mathrm{A})
\]

is bijective: indeed, it suffices to verify that, for every S \(^{-1}\) A-module M, every k-derivation D : A → M extends uniquely to a k-derivation D : S \(^{-1}\) A → M; this follows from the argument of loc. cit., p. 123, Proposition 5.

Let \(k\) be a field and A a local, essentially of finite type \(k\)-algebra. Set \(n = \dim(A) + \deg.\mathrm{tr}_k(\kappa_A)\). Let B be a finitely generated \(k\)-algebra and \(\mathfrak{q}\) a prime ideal of B such that the \(k\)-algebra A is isomorphic to \(\mathbf{B}_{\mathfrak{q}}\) (\(cf.\ n^{\circ} 1\), Example 2); we have \(n = \dim_{\mathfrak{q}}(\mathbf{B})\) (VIII, § 2, no. 4, Corollary 5 of Theorem 3). By loc. cit., Theorem 3, c) and Corollary 5, we also have

\[
n = \sup \deg . \operatorname{tr} _ {k} (\kappa (\mathfrak {p}))
\]

where p ranges over the finite family of minimal prime ideals of A. If A is integral, n is therefore the transcendence degree of the field of fractions of A.

THEOREM 1. Let \(k\) be a field and A a local, essentially of finite type \(k\)-algebra. Set \(n = \dim(A) + \deg.\mathrm{tr}_{k}(\kappa_{A})\).

\begin{enumerate}
\def\labelenumi{\alph{enumi})}
\item
  One has \([\kappa_{A} \otimes_{A} \Omega_k(A): \kappa_{A}] \geqslant n\).
\item
  The following conditions are equivalent:
\end{enumerate}

\begin{enumerate}
\def\labelenumi{(\roman{enumi})}
\item
  the k-algebra A is absolutely regular;
\item
  the A-module \(\Omega_{k}(\mathrm{A})\) is free of rank \(n\);
\end{enumerate}

\[
\left[ \kappa_ {A} \otimes_ {A} \Omega_ {k} (A): \kappa_ {A} \right] = n
\]

Consider the assertion

(iii') one has \([\kappa_{\mathrm{A}} \otimes_{A} \Omega_k(\mathrm{A}):\kappa_{\mathrm{A}}] \leqslant n\);

it is clear that (ii) implies (iii) and that (iii) implies (iii'). To prove the theorem it suffices to prove the implications (i) \(\Rightarrow\) (ii) and (iii') \(\Rightarrow\) (i).

Choose a finitely generated \(k\)-algebra B and a prime ideal \(\mathfrak{q}\) of B such that A is isomorphic to \(\mathrm{B}_{\mathfrak{q}}\) (\(n^{\circ} 1\), Example 2). We have \(\dim_{\mathfrak{q}}(\mathrm{B}) = n\); replacing B by a ring \(\mathrm{B}_f\), with \(f \in \mathrm{B} - \mathfrak{q}\), one can impose that the spectrum of B is connected and of dimension \(n\), which we will henceforth assume.

\begin{enumerate}
\def\labelenumi{(\roman{enumi})}
\tightlist
\item
  \(\Rightarrow\) (ii): Suppose the \(k\)-algebra A is absolutely regular. As in the proof of Proposition 8, denote by \(\mu : A \otimes_k A \to A\) and \(\nu : \kappa_A \otimes_k A \to \kappa_A\) the homomorphisms deduced from the multiplication of A; set I = \(\operatorname{Ker}(\mu)\) and \(\mathfrak{n} = \operatorname{Ker}(\nu)\). The ideal I is completely secant; the ideal \(\mathfrak{n}\) is maximal and the local ring (\(\kappa_A \otimes_k A\) ) \(_{\mathfrak{n}}\) is regular (loc. cit.). By Theorem 1 of § 5, No. 2, the \(A\)-module of finite type I/I², which is none other than \(\Omega_k(A)\), is projective, hence free. But according to the remark of § 5, No. 6, the ideal \(\mathfrak{n}\) is identified with \(\kappa_A \otimes_A 1\), so that \(\mathfrak{n}/\mathfrak{n}^2\) is isomorphic to \(\kappa_A \otimes_A \Omega_k(A)\). One therefore has
\end{enumerate}

\[
\operatorname{rg} \left(\Omega_ {k} (\mathrm{A})\right) = \left[ \kappa_ {\mathrm{A}} \otimes_ {k} \Omega_ {k} (\mathrm{A}): \kappa_ {\mathrm{A}} \right] = \left[ \mathfrak {n} / \mathfrak {n} ^ {2}: \kappa_ {\mathrm{A}} \right] = \dim \left(\kappa_ {\mathrm{A}} \otimes_ {k} \mathrm{A}\right) _ {\mathfrak {n}}.
\]

Let B' denote the ring \(\kappa_{A}\otimes_{k}B\) , and let S be the image in B' of B-q. The ring \(\kappa_{A}\otimes_{k}A\) is identified with B' \(\otimes_{B}B_{q}\) , that is, with S \(^{-1}\) B'. There therefore exists a maximal ideal \(q'\) of \(B'\) not meeting S such that one has \(n = S^{-1}q'\) ; the ring \((\kappa_A \otimes_k A)_n\) is then isomorphic to \(B_{q'}'\) and one has

\[
\dim (\kappa_ {A} \otimes_ {k} A) _ {\mathfrak {n}} = \dim (\kappa_ {A} \otimes_ {k} B) _ {\mathfrak {q} ^ {\prime}} = \dim_ {\mathfrak {q} ^ {\prime}} (\kappa_ {A} \otimes_ {k} B).
\]

Since n contains \(\kappa_{\mathrm{A}} \otimes_{k} qB_{q}\) , \(q'\) contains the ideal \(\kappa_{\mathrm{A}} \otimes q\) of \(B'\) , hence its inverse image in B contains \(q\) ; it is equal to \(q\) since \(q'\) does not meet S. By VIII, § 2, no. 4, cor. of prop. 5, we have

\[
\dim_ {\mathfrak {q} ^ {\prime}} \left(\kappa_ {A} \otimes_ {k} B\right) = \dim_ {\mathfrak {q}} (B) = n,
\]

which proves (ii).

(iii') \(\Rightarrow\) (i): Suppose that we have \([\kappa_{\mathrm{A}}\otimes_{\mathrm{A}}\Omega_{k}(A):\kappa_{\mathrm{A}}]\leqslant n\), that is \([\kappa (\mathfrak{q})\otimes_{\mathrm{B}}\Omega_{k}(\mathrm{B}):\kappa (\mathfrak{q})]\leqslant n\), and let us prove that the \(k\)-algebra B is absolutely regular. Let \((x_{1},\ldots ,x_{n})\) be a sequence of elements of B such that \(1\otimes dx_1,\dots ,1\otimes dx_n\) generate the \(\kappa (\mathfrak{q})\)-vector space \(\kappa (\mathfrak{q})\otimes_{\mathrm{B}}\Omega_{k}(\mathrm{B})\). Replacing B by \(\mathrm{B}_f\), for a suitable element \(f\) of \(\mathrm{B} - \mathfrak{q}\), one may assume that \(dx_{1},\ldots ,dx_{n}\) generate the B-module \(\Omega_k(\mathrm{B})\). By Nakayama's lemma and II, § 5, No. 1, Proposition 2, let \(\bar{k}\) be an algebraically closed extension of \(k\). It suffices for us to prove that the \(\bar{k}\)-algebra \(\mathrm{B}_{(\bar{k})}\) is regular, since this will imply that B is absolutely regular (Corollary 4 of Proposition 7 of No. 4). For every canonical component C of \(\mathrm{B}_{(k)}\), the differentials \(d(1\otimes x_i)\) generate the C-module \(\Omega_{\bar{k}}(\mathrm{C})\) (A, III, § 10, No. 12, Proposition 20). The theorem therefore follows from the following two lemmas:

Lemma 4.—Let \(k\) be a field, and let K be an algebraic extension of \(k\). Let B be a \(k\)-algebra whose spectrum is connected. For every canonical component C of \(\mathrm{B}_{(\mathrm{K})}\), one has \(\dim(\mathrm{C}) = \dim(\mathrm{B})\).

We first treat the case where the extension K is of finite degree over \(k\). The B-module C is a direct summand of a free module, hence finitely generated projective. Since the function \(\mathfrak{p} \mapsto \mathrm{rg}_{\mathfrak{p}}\) is constant on \(\operatorname{Spec}(B)\) (II, § 5, no. 3), the support of the B-module C is equal to \(\operatorname{Spec}(B)\), so that the B-module C is faithfully flat. This implies \(\dim(C) \geqslant \dim(B)\) (VIII, § 2, no. 1, prop. 2), whence the desired equality, since we have \(\dim(C) \leqslant \dim(B_{(K)}) = \dim(B)\) (VIII, § 2, no. 4, cor. of prop. 5).

For the general case, let \(e\) be the idempotent element of \(B_{(K)}\) such that \(C = B_{(K)}/B_{(K)}e\). There exists a subextension \(K'\) of \(K\) of finite degree such that \(e\) is the image of an idempotent element \(e'\) of \(B_{(K')}\). Set \(C' = B_{(K')} / B_{(K')}e'\). The ring \(C' \otimes_{K'} K\) is identified with \(C\), and \(C'\) is a canonical component of \(B_{(K')}\). One has \(\dim(C') = \dim(B)\) by the case already treated, and \(\dim(C') = \dim(C)\) by loc. cit., whence the lemma.

Lemma 5.—Let \(k\) be an algebraically closed field, and let A be a finite-type \(k\)-algebra whose spectrum is connected, let \(n\) be an integer, and let \((x_{1},\ldots,x_{n})\) be a finite sequence of elements of A. We assume that A has dimension \(n\) and that the differentials \(dx_{1},\ldots,dx_{n}\) generate the A-module \(\Omega_{k}(\mathrm{A})\). Then the ring A is integral and regular, and the \(A\)-module \(\Omega_{k}(\mathrm{A})\) is free with basis \((dx_{1},\ldots,dx_{n})\).

Let \(\mathfrak{m}\) be a maximal ideal of A such that \(\dim(A_{\mathfrak{m}})=n\). One has \([A/\mathfrak{m}:k]=1\) (V, § 3, No. 3, Proposition 1 (iii)), hence \(A=\mathfrak{m}\oplus k1_{A}\). Let \(p\) and \(q\) be the corresponding projectors

For \(a\) and \(b\) in \(A\), one has

\[
a b = \left(p (a) q (b) + q (a) p (b) + p (a) p (b), q (a) q (b)\right),
\]

whence \(p(ab) \equiv p(a)q(b) + q(a)p(b) \pmod{\mathfrak{m}^{2}}\). Consequently, the map \(\delta: A \to \mathfrak{m}/\mathfrak{m}^{2}\) which associates to each element x of A the class of \(p(x)\) modulo \(\mathfrak{m}^{2}\) is a k-derivation of A into the k-vector space \(\mathfrak{m}/\mathfrak{m}^{2}\). There therefore exists an A-linear map \(\phi: \Omega_{k}(A) \to \mathfrak{m}/\mathfrak{m}^{2}\) such that \(\delta(x) = \phi(dx)\) for every \(x \in A\). Since \(\delta\) is surjective, the \(\phi(dx_{i})\) generate the A/m-vector space \(\mathfrak{m}/\mathfrak{m}^{2}\), and one has \([\mathfrak{m}/\mathfrak{m}^{2}: A/\mathfrak{m}] \leqslant n = \dim(A_{\mathfrak{m}})\). It follows that \(A_{\mathfrak{m}}\) is regular and that the images of \(dx_{i}\) form a basis of the A/m-vector space A/m \(\otimes_{A} \Omega_{k}(A)\).

Now let us prove that the ring A is integral and regular. There exists a minimal prime ideal q of A such that \(\dim(A/q) = n\). For every maximal ideal m of A containing q, we have \(\dim(A_{\mathfrak{m}}) = n\) (VIII, § 2, No. 4, Corollary 2 of Theorem 3), so that \(A_{\mathfrak{m}}\) is regular by what we have just seen. In particular, \(A_{\mathfrak{m}}\) is integral, which implies \(qA_{\mathfrak{m}} = 0\). Since this holds for all maximal ideals m of V(q), we deduce that \(\operatorname{Supp}(\mathfrak{q}) \cap V(\mathfrak{q}) = \varnothing\). But \(\operatorname{Spec}(A)\) is connected, V(q) is nonempty, and we have \(\operatorname{Supp}(\mathfrak{q}) \cup V(\mathfrak{q}) = \operatorname{Spec}(A)\) (II, § 4, No. 4, Proposition 16). From this we deduce \(\operatorname{Supp}(\mathfrak{q}) = \varnothing\), whence q = 0, which means that A is an integral domain. We then have \(\dim(A_{\mathfrak{m}}) = n\) for every maximal ideal m of A; applying the first part of the proof, we deduce that A is regular.

Finally, let \(\sum_{i=1}^{n} a_i dx_i = 0\) be a linear relation among the \(dx_i\) with coefficients in \(A\). If the \(a_i\) are not all zero, there exists an index \(i\) and a maximal ideal \(\mathfrak{m}\) of A such that \(a_i\) does not belong to \(\mathfrak{m}\) (V, § 3, No. 3, Proposition 1, (iii) and (iv)); but this contradicts the fact proved above that the classes of the \(dx_i\) in (A/m) \(\otimes_{\mathrm{A}} \Omega_k(\mathrm{A})\) are linearly independent.

Example.--- When A is a finitely generated extension of k, Theorem 1 recovers Corollary 1 of A, V, p. 128, in view of Example 2 of No. 4.

COROLLARY 1. - Let \(k\) be a field and \(A\) an essentially of finite type \(k\)-algebra. The set of prime ideals \(\mathfrak{p}\) of \(\operatorname{Spec}(A)\) such that the \(k\)-algebra \(\mathrm{A}_{\mathfrak{p}}\) is absolutely regular is open in \(\operatorname{Spec}(A)\).

We may assume that the \(k\)-algebra A is finitely generated. The set under consideration then consists of the prime ideals \(\mathfrak{p}\) such that \([\kappa(\mathfrak{p}) \otimes_k \Omega_k(A) : \kappa(\mathfrak{p})] \leqslant \dim_{\mathfrak{p}}(A)\). Now the function \(\mathfrak{p} \mapsto \dim_{\mathfrak{p}}(A)\) is lower semicontinuous by definition, and the function \(\mathfrak{p} \mapsto [\kappa(\mathfrak{p}) \otimes_k \Omega_k(A) : \kappa(\mathfrak{p})]\) is upper semicontinuous (by Nakayama's lemma and II, § 5, No. 1, Proposition 2).

We shall see later (§ 7, No. 9, Corollary 4 of Theorem 3) that under the hypotheses of Corollary 1, the set of prime ideals p of A such that the ring \(A_{p}\) is regular is open in \(\operatorname{Spec}(A)\).

COROLLARY 2.--- Let \(k\) be a field and A an essentially finitely generated \(k\)-algebra. For A to be absolutely regular, it is necessary and sufficient that the A-module \(\Omega_{k}(A)\) be projective and that, for every minimal prime ideal \(\mathfrak{q}\) of A, the \(k\)-algebra \(\mathrm{A}_{\mathfrak{q}}\) be separable.

Suppose A is absolutely regular. For every prime ideal p of A, the k-algebra \(A_{p}\) is absolutely regular, so that the \(A_{p}\)-module \(\Omega_{k}(A_{p})\) is free (Theorem 1); it follows that the A-module \(\Omega_{k}(A)\) is projective (II, § 5, No. 2, Theorem 1). Moreover, for every minimal prime ideal q of A, the local k-algebra \(A_{q}\) is Artinian and absolutely regular, hence is a field, a separable extension of k (No. 4, Example 2).

Conversely, suppose that the A-module \(\Omega_{k}(\mathbf{A})\) is projective and that the \(k\)-algebra \(\mathbf{A}_{\mathfrak{q}}\) is separable for every minimal prime ideal \(\mathfrak{q}\) of \(\mathbf{A}\). Let \(\mathfrak{p}\) be a prime ideal of \(\mathbf{A}\), and let \(\mathfrak{q}\) be a minimal prime ideal of \(\mathbf{A}\) contained in \(\mathfrak{p}\). Since the \(\mathbf{A}_{\mathfrak{p}}\)-module \(\Omega_{k}(\mathbf{A})_{\mathfrak{p}}\) is free (II, § 3, No. 2, Corollary 2 of Proposition 5), we have

\[
[ \kappa (\mathfrak {p}) \otimes_ {\mathrm{A}} \Omega_ {k} (\mathrm{A}): \kappa (\mathfrak {p}) ] = [ \kappa (\mathfrak {q}) \otimes_ {\mathrm{A}} \Omega_ {k} (\mathrm{A}): \kappa (\mathfrak {q}) ].
\]

The \(k\)-algebra \(\mathrm{A}_{\mathfrak{q}}\) is Artinian and separable, hence absolutely regular (\(n^{\circ}\) 4, Example 2). Theorem 1 implies

\[
[ \kappa (\mathfrak {q}) \otimes_ {\mathrm{A}} \Omega_ {k} (\mathrm{A}): \kappa (\mathfrak {q}) ] = \deg . \operatorname{tr} _ {k} (\kappa (\mathfrak {q})).
\]

It then follows from Theorem 1 that the \(k\)-algebra \(\mathrm{A}_{\mathfrak{p}}\) is absolutely regular, whence the corollary.

Note.--- Suppose the k-algebra A is absolutely regular. Then the total ring of fractions F of A is identified with the product of the \(A_{q}\), where q ranges over the set of minimal prime ideals of A (IV, § 2, No. 5, Proposition 10); it is therefore a separable k-algebra. Suppose conversely that F is a separable k-algebra; for every minimal prime ideal q of A, the k-algebra \(A_{q}\) is a ring of fractions of F (IV, § 1, No. 1, Corollary 3 of Proposition 2 and No. 3, Corollary 1 of Proposition 7), hence a separable k-algebra (No. 4, Proposition 6). If moreover the A-module \(\Omega_{k}(A)\) is projective, it follows from Corollary 2 that the k-algebra A is absolutely regular.

COROLLARY 3.--- Let \(k\) be a field of characteristic 0 and A an essentially of finite type \(k\)-algebra. For A to be regular, it is necessary and sufficient that the A-module \(\Omega_{k}(\mathrm{A})\) be projective.

By Corollary 2, it suffices to prove that a local Artinian \(k\)-algebra A such that \(\Omega_k(A)\) is a free A-module is a field. Now, by IX, § 3, No. 3, Theorem 1, there exists a subfield K of A such that \(A = K \oplus \mathfrak{m}_A\). Let \(\delta: A \to \mathfrak{m}_A / \mathfrak{m}_A^2\) be the composite of the projection of A onto \(\mathfrak{m}_A\) and the canonical map \(\mathfrak{m}_A \to \mathfrak{m}_A / \mathfrak{m}_A^2\). One verifies as in the proof of Lemma 4 that \(\delta\) is a k-derivation. But every \(k\)-derivation of A into an A-module annihilates \(\mathfrak{m}_A\): indeed, let \(x \in \mathfrak{m}_A\), and choose an integer \(n \geqslant 1\) such that \(x^{n-1} \neq 0\) and \(x^n = 0\). This implies, in the A-module \(\Omega_k(A)\), the relation \(nx^{n-1} dx = 0\), hence \(dx = 0\) since \(nx^{n-1} \neq 0\). Thus \(\delta(\mathfrak{m}_A) = 0\), hence \(\mathfrak{m}_A = \mathfrak{m}_A^2\), and ultimately \(\mathfrak{m}_A = \{0\}\).

